\chapter{Entraînement et prédiction}
\label{ch:implementation}

\begin{lecture}
Deux programmes suffisent à produire le classifieur. Le premier lit les élèves
étiquetés et écrit les paramètres ; le second lit ces paramètres et attribue une
maison à chaque nouvel élève. Ce chapitre donne les colonnes qui entrent dans le
calcul, les dimensions de chaque tableau, les valeurs des hyperparamètres et le
contenu des fichiers échangés. Les nombres cités proviennent de
\texttt{dataset\_train.csv}.
\end{lecture}

\section{Colonnes du fichier}
\label{sec:features}

Le fichier d'apprentissage contient $1\,600$ lignes et dix-neuf colonnes.

\begin{table}[htbp]
\centering
\begin{tabularx}{\textwidth}{>{\ttfamily\footnotesize}l X}
\toprule
\textrm{colonne} & \textrm{contenu et emploi} \\
\midrule
Index & entier de $0$ à $1\,599$. Recopié dans \texttt{houses.csv} pour apparier chaque prédiction à son élève \\
Hogwarts House & la réponse $c_i$. Quatre valeurs, d'effectifs $529$, $443$, $327$ et $301$. Source des cibles binaires $y_{ik}$ \\
First Name, Last Name & identité de l'élève. Écartées \\
Birthday & date de naissance au format \texttt{AAAA-MM-JJ}, de $1996$ à $2001$. Écartée d'après la section~\ref{sec:hand-birthday} \\
Best Hand & \texttt{Left} pour $790$ élèves, \texttt{Right} pour $810$. Écartée d'après la section~\ref{sec:hand-birthday} \\
\textrm{treize colonnes de notes} & nombres réels. Variables candidates \\
\bottomrule
\end{tabularx}
\caption{Composition du fichier d'apprentissage.}
\label{tab:columns}
\end{table}

\section{Best Hand et Birthday}
\label{sec:hand-birthday}

Ces deux colonnes sont remplies sur les $1\,600$ lignes et demandent un codage
numérique avant d'entrer dans le modèle. Trois mesures décident de leur sort.

La proportion de gauchers varie de $48{,}1\,\%$ chez Ravenclaw à $51{,}4\,\%$
chez Gryffindor, pour $49{,}4\,\%$ sur l'ensemble. Le tableau de contingence
$4\times2$ donne un \rsym{chi2}{$\chi^2$} de $0{,}94$ à trois degrés de liberté, quand le
seuil à $5\,\%$ vaut $7{,}81$.

L'écart entre la note moyenne des gauchers et celle des droitiers reste sous
$0{,}10$ écart-type pour les treize cours, avec un maximum de $0{,}095$ pour
Defense Against the Dark Arts. L'\rsym{erreurtype}{erreur type} d'un tel écart vaut
$\sqrt{1/790+1/810}=0{,}050$ écart-type : les treize valeurs tiennent à moins de
deux erreurs types de zéro. Les maisons, elles, se séparent de $1{,}86$ à
$2{,}66$ écarts-types sur ces mêmes cours.

Le mois de naissance donne un $\chi^2$ de $33{,}4$ à $33$ degrés de liberté.
L'année donne $28{,}1$ à $15$ degrés de liberté, au-dessus du seuil à $5\,\%$ qui
vaut $25{,}0$ : une association faible, portée en partie par les $40$ élèves nés
en $1996$ contre environ $320$ pour chacune des années suivantes.

Le test décisif porte sur l'exactitude. Best Hand codée $\pm1$, puis l'année de
naissance standardisée, ajoutées chacune comme onzième variable, laissent
l'exactitude hors échantillon à $0{,}9832$ sur cinq graines de découpage,
identique aux quatre décimales. Les coefficients de Best Hand s'échelonnent de
$-0{,}20$ à $+0{,}19$, contre $2{,}39$ pour Astronomy chez Hufflepuff. Les deux
colonnes restent hors du modèle.

\section{Redondance}

\begin{lecture}
Deux colonnes qui portent la même information gênent le calcul. Le modèle peut
augmenter le coefficient de l'une et diminuer celui de l'autre sans rien changer
à ses prédictions : ses coefficients dérivent alors sans jamais se fixer. Une des
deux colonnes suffit.
\end{lecture}

Le jeu de données vérifie
\[
  \text{Astronomy}=-100\times\text{Defense Against the Dark Arts}
\]
sur les $1\,537$ lignes où les deux valeurs sont présentes, avec un écart maximal
de $4{,}5\cdot10^{-13}$ et une \rsym{correlation}{corrélation} de $-1$ à six décimales. Le
chapitre~\ref{ch:likelihood} établit la conséquence : deux colonnes
\rsym{colinearite}{proportionnelles} rendent \rsym{xtwx}{$\X^\top W\X$} singulière et laissent la fonction de coût
plate le long d'une direction. Les coefficients dérivent alors librement sur
cette direction pendant que les scores restent identiques. Astronomy est
conservée, Defense Against the Dark Arts sert à l'imputation décrite plus bas.

\section{Pouvoir discriminant}

\begin{lecture}
Une note aide à reconnaître une maison lorsque les quatre maisons y obtiennent
des résultats différents. Si les quatre ont la même moyenne, la note ne distingue
personne. Le nombre $\delta_j$ ci-dessous mesure cet écart entre maisons, dans
une unité qui permet de comparer des cours notés sur des échelles très
différentes.
\end{lecture}

Pour la variable $j$, l'étendue des moyennes par maison rapportée à l'écart-type
global vaut
\[
  \delta_j=\frac{\max_k\bar r_{jk}-\min_k\bar r_{jk}}{s_j}.
\]
Un $\delta_j$ de $2$ signifie que les maisons extrêmes ont des moyennes séparées
par deux écarts-types : leurs distributions se chevauchent peu et la variable
sépare. Un $\delta_j$ de $0{,}07$ place les quatre moyennes à l'intérieur d'un
quatorzième d'écart-type : les quatre histogrammes se superposent.

\begin{table}[htbp]
\centering
\setlength{\tabcolsep}{7pt}
\begin{tabular}{lrrl}
\toprule
variable & $\delta_j$ & manquantes & décision \\
\midrule
\cellcolor{warning!14}Arithmancy                & \cellcolor{warning!14}$0{,}068$ & \cellcolor{warning!14}$34$ & \cellcolor{warning!14}écartée : homogène \\
\cellcolor{warning!14}Care of Magical Creatures & \cellcolor{warning!14}$0{,}147$ & \cellcolor{warning!14}$40$ & \cellcolor{warning!14}écartée : homogène \\
\midrule
Ancient Runes                & $1{,}860$ & $35$ & retenue \\
Herbology                    & $1{,}879$ & $33$ & retenue \\
\cellcolor{warning!14}Defense Against the Dark Arts & \cellcolor{warning!14}$1{,}909$ & \cellcolor{warning!14}$31$ & \cellcolor{warning!14}écartée : redondante \\
Astronomy                    & $1{,}911$ & $32$ & retenue \\
Muggle Studies               & $2{,}038$ & $35$ & retenue \\
Potions                      & $2{,}076$ & $30$ & retenue \\
History of Magic             & $2{,}221$ & $43$ & retenue \\
Transfiguration              & $2{,}284$ & $34$ & retenue \\
Divination                   & $2{,}368$ & $39$ & retenue \\
Charms                       & $2{,}466$ & $0$  & retenue \\
Flying                       & $2{,}657$ & $0$  & retenue \\
\bottomrule
\end{tabular}
\caption{Les treize notes classées par pouvoir discriminant croissant. Arithmancy
et Care of Magical Creatures sont en dessous des autres d'un facteur douze ; la
coupure tombe dans un intervalle vide.}
\label{tab:feature-selection}
\end{table}

\begin{resultat}[Variables retenues]
Astronomy, Herbology, Divination, Muggle Studies, Ancient Runes,
History of Magic, Transfiguration, Potions, Charms, Flying.
Donc $n=10$, $\X\in\R^{m\times11}$ et chaque $\thv_k\in\R^{11}$.
\end{resultat}

\section{Valeurs manquantes}

Onze des treize colonnes de notes comportent des cases vides, entre $30$ et $43$
chacune, réparties sur $349$ lignes distinctes.

Les $1\,600$ lignes possèdent au moins une des deux valeurs Astronomy ou Defense
Against the Dark Arts. La relation de proportionnalité donne donc les $32$
valeurs manquantes d'Astronomy à partir de la colonne écartée, au facteur $-100$
près. Ces $32$ cases reçoivent leur valeur exacte, et Astronomy devient complète.

Restent alors $239$ lignes portant au moins une case vide parmi les dix variables
retenues, soit $14{,}9\,\%$ de l'échantillon : les supprimer ramènerait
l'apprentissage à $1\,361$ élèves. L'imputation médiane du
chapitre~\ref{ch:data} les conserve toutes, et porte sur les sept colonnes encore
incomplètes. Charms, Flying et Astronomy sont pleines.

\section{Matrice de design}

Le programme applique dans l'ordre : reconstruction d'Astronomy, imputation
médiane, standardisation, ajout de la colonne de biais.

\begin{procedure}{Construire $\X$}
\textbf{Entrée :} un fichier de données. \textbf{Sorties :} $\X$, le vecteur des
identifiants \texttt{Index}, et le vecteur des classes lorsque la colonne est
remplie.
\begin{enumerate}
  \item lire les lignes dans l'ordre du fichier ; cet ordre est celui de
        \texttt{houses.csv} ;
  \item extraire les dix colonnes retenues, dans l'ordre fixé par le fichier de
        poids ;
  \item convertir les cases en nombres réels ; une case vide devient une valeur
        manquante ;
  \item remplir les valeurs manquantes d'Astronomy par $-100\,r_{\mathrm{DADA}}$,
        puis les valeurs manquantes restantes par la \rsym{mediane}{médiane} d'apprentissage de
        leur colonne ;
  \item appliquer $x_{ij}=(r_{ij}-\mu_j)/s_j$ avec les $\mu_j$ et $s_j$
        d'apprentissage ;
  \item préfixer chaque ligne d'une coordonnée égale à $1$, ce qui donne
        $\X\in\R^{m\times11}$.
\end{enumerate}
\end{procedure}

Sur \texttt{dataset\_train.csv}, $m=1\,600$ et $\X$ mesure $1\,600\times11$. Sur
\texttt{dataset\_test.csv}, $m=400$ et $\X$ mesure $400\times11$. Le découpage
stratifié $80\,\%/20\,\%$ du chapitre~\ref{ch:evaluation} partage les $1\,600$
lignes d'apprentissage en $1\,278$ et $322$.

\section{Le programme d'entraînement}

\begin{lecture}
L'entraînement répète quatre calculs jusqu'à ce que les coefficients cessent de
bouger : un score pour chaque élève, la probabilité qui lui correspond, l'écart
entre cette probabilité et la bonne réponse, puis une correction des coefficients
proportionnelle à cet écart. La même boucle tourne quatre fois, une par maison.
\end{lecture}

L'entraînement traite les quatre maisons l'une après l'autre. Pour la maison $k$,
il forme la cible $\y_k\in\{0,1\}^{1600}$ par $y_{ik}=\ind[c_i=k]$, part de
$\thv_k=\mathbf 0$, puis répète quatre opérations :
\[
  \mathbf z=\X\thv_k\;(m),\qquad
  \mathbf p=\sigma(\mathbf z)\;(m),\qquad
  \mathbf g=\tfrac1m\X^\top(\mathbf p-\y_k)\;(11),\qquad
  \thv_k\leftarrow\thv_k-\alpha\mathbf g.
\]
Le coût $J$ de l'équation~\eqref{eq:loss} sert au critère d'arrêt, le gradient
$\mathbf g$ est l'équation~\eqref{eq:gradient}. Le départ à $\thv_k=\mathbf 0$
donne un résultat reproductible d'une exécution à l'autre, la convexité établie
au chapitre~\ref{ch:likelihood} garantissant l'unicité de la limite à direction
plate près.

\begin{procedure}{logreg\_train}
\textbf{Entrée :} \texttt{dataset\_train.csv}. \textbf{Sortie :} le fichier de
poids.
\begin{enumerate}
  \item calculer médianes, $\mu_j$ et $s_j$ sur les lignes d'apprentissage, puis
        construire $\X$ ;
  \item pour $k=1,\ldots,4$, former $\y_k$ et itérer les quatre opérations
        ci-dessus jusqu'à $|J^{(t+1)}-J^{(t)}|/\max(1,|J^{(t)}|)<\eps$ ou
        épuisement du plafond d'itérations ;
  \item relever pour chaque maison le coût initial, le coût final, le nombre
        d'itérations et $\|\mathbf g\|_2$ ;
  \item écrire le fichier de poids avec les statistiques de préparation.
\end{enumerate}
\end{procedure}

\begin{table}[htbp]
\centering
\begin{tabular}{rrrrr}
\toprule
$\alpha$ & $\eps$ & itérations & $\|\nabla J\|_2$ final & exactitude hors échantillon \\
\midrule
$0{,}5$ & $10^{-6}$ & $1\,817$ & $1{,}4\cdot10^{-3}$ & $0{,}981$ \\
$1$     & $10^{-6}$ & $1\,461$ & $1{,}0\cdot10^{-3}$ & $0{,}981$ \\
$2$     & $10^{-6}$ & $1\,150$ & $7{,}1\cdot10^{-4}$ & $0{,}981$ \\
$1$     & $10^{-7}$ & $5\,448$ & $3{,}2\cdot10^{-4}$ & $0{,}981$ \\
\bottomrule
\end{tabular}
\caption{Itérations consommées par la maison la plus lente, sur les $1\,278$
lignes d'apprentissage du découpage stratifié. Diviser $\eps$ par dix multiplie
les itérations par près de quatre et laisse les $322$ décisions du test interne
inchangées : la convergence numérique se poursuit après que le classement s'est
stabilisé.}
\label{tab:training-defaults}
\end{table}

\begin{resultat}[Configuration de départ]
$\alpha=1$, $\eps=10^{-6}$, plafond de $6\,000$ itérations, $\thv_k=\mathbf 0$.
Cette configuration atteint $0{,}981$ sur un test interne de $322$ élèves, et
$0{,}983$ en moyenne sur cinq graines de découpage.
\end{resultat}

\section{Le fichier de poids}
\label{sec:weights}

La prédiction a besoin de trois choses : l'ordre des dix variables, les
statistiques qui les transforment, et les quatre vecteurs de paramètres associés
à leurs libellés. Les statistiques figurent dans le fichier parce qu'elles sont
estimées sur l'apprentissage ; les recalculer sur \texttt{dataset\_test.csv}
placerait les scores dans une autre échelle et ferait chuter l'exactitude.

\begin{fichier}[weights.csv]
feature,median,mu,sigma\\
Astronomy,258.1,38.6,520.8\\
Herbology,3.5,1.2,5.2\\
{[}...{]}\\
Flying,-2.5,22.0,97.6\\
class,bias,Astronomy,Herbology,...,Flying\\
Gryffindor,-3.54,0.50,-1.81,...,1.61\\
Hufflepuff,-2.36,2.48,1.63,...,-0.48\\
Ravenclaw,-2.56,-1.37,0.74,...,-0.17\\
Slytherin,-3.77,-2.02,-0.70,...,-0.59
\end{fichier}

Ces valeurs sont celles produites par le programme de la
section~\ref{sec:pseudo-train} sur les $1\,600$ lignes du fichier
d'apprentissage. L'écart-type d'Astronomy vaut $521$ et celui d'Herbology
$5{,}2$ : un facteur cent entre deux colonnes que le même $\alpha$ doit faire
progresser, ce que la standardisation ramène à $1$.

L'ordre des colonnes et l'ordre des classes appartiennent au modèle. Un vecteur
$\thv_k$ s'applique aux variables dans l'ordre où il a été ajusté, et l'indice
renvoyé par $\Argmax$ se lit dans la liste des libellés enregistrée à côté de
lui. Les deux ordres figurent explicitement dans le fichier.

\section{Primitives à écrire}

Le sujet écarte les fonctions toutes faites de statistique descriptive. Cinq
primitives suffisent, toutes en une passe ou deux sur un tableau.

\begin{pseudo}[primitives]
fonction moyenne(v) :\\
\hspace*{2em}s $\leftarrow$ 0\\
\hspace*{2em}pour x dans v : s $\leftarrow$ s + x\\
\hspace*{2em}retourner s / longueur(v)\\
\\
fonction ecart\_type(v) :\\
\hspace*{2em}m $\leftarrow$ moyenne(v) ; s $\leftarrow$ 0\\
\hspace*{2em}pour x dans v : s $\leftarrow$ s + (x - m)$^2$\\
\hspace*{2em}retourner racine(s / longueur(v))\\
\\
fonction mediane(v) :\\
\hspace*{2em}t $\leftarrow$ trier(v) ; n $\leftarrow$ longueur(t)\\
\hspace*{2em}si n impair : retourner t[n // 2]\\
\hspace*{2em}retourner (t[n // 2 - 1] + t[n // 2]) / 2\\
\\
fonction sigmoide(z) :\\
\hspace*{2em}si z $\geq$ 0 : retourner 1 / (1 + exp(-z))\\
\hspace*{2em}retourner exp(z) / (1 + exp(z))\\
\\
fonction softplus(z) :\\
\hspace*{2em}retourner max(z, 0) + log(1 + exp(-abs(z)))
\end{pseudo}

Les deux dernières viennent du chapitre~\ref{ch:optim} : leur exponentielle porte
toujours sur un nombre négatif ou nul, ce qui les garde finies pour
$z=\pm1000$.

\section{Pseudo-code de logreg\_train}
\label{sec:pseudo-train}

Le tableau \texttt{X} a $m$ lignes et $11$ colonnes, \texttt{theta[k]} $11$
composantes, \texttt{y} $m$ composantes. Les indices de colonne vont de $0$ à
$10$, la colonne $0$ portant le biais.

\begin{pseudo}[logreg\_train dataset\_train.csv]
VARIABLES $\leftarrow$ [Astronomy, Herbology, Divination, Muggle Studies,\\
\hspace*{5em}Ancient Runes, History of Magic, Transfiguration,\\
\hspace*{5em}Potions, Charms, Flying]\\
MAISONS $\leftarrow$ [Gryffindor, Hufflepuff, Ravenclaw, Slytherin]\\
ALPHA $\leftarrow$ 1.0 ; EPS $\leftarrow$ 1e-6 ; MAX\_ITER $\leftarrow$ 6000\\
\\
lignes $\leftarrow$ lire\_csv(argv[1])\\
\\
pour i dans lignes :\\
\hspace*{2em}si lignes[i][Astronomy] est vide et lignes[i][DADA] est rempli :\\
\hspace*{4em}lignes[i][Astronomy] $\leftarrow$ -100 * lignes[i][DADA]\\
\\
pour f dans VARIABLES :\\
\hspace*{2em}presentes $\leftarrow$ [lignes[i][f] pour i tel que lignes[i][f] rempli]\\
\hspace*{2em}med[f] $\leftarrow$ mediane(presentes)\\
\hspace*{2em}pour i dans lignes :\\
\hspace*{4em}si lignes[i][f] est vide : lignes[i][f] $\leftarrow$ med[f]\\
\hspace*{2em}mu[f] $\leftarrow$ moyenne(colonne f)\\
\hspace*{2em}sd[f] $\leftarrow$ ecart\_type(colonne f)\\
\hspace*{2em}si sd[f] = 0 : erreur("colonne constante", f)\\
\\
X $\leftarrow$ tableau[m][11]\\
pour i de 0 a m-1 :\\
\hspace*{2em}X[i][0] $\leftarrow$ 1.0\\
\hspace*{2em}pour j de 0 a 9 :\\
\hspace*{4em}f $\leftarrow$ VARIABLES[j]\\
\hspace*{4em}X[i][j+1] $\leftarrow$ (lignes[i][f] - mu[f]) / sd[f]\\
\\
pour k de 0 a 3 :\\
\hspace*{2em}y $\leftarrow$ [1 si lignes[i][House] = MAISONS[k] sinon 0]\\
\hspace*{2em}th $\leftarrow$ [0.0] * 11\\
\hspace*{2em}J\_prec $\leftarrow$ indefini\\
\hspace*{2em}pour t de 1 a MAX\_ITER :\\
\hspace*{4em}J $\leftarrow$ 0 ; g $\leftarrow$ [0.0] * 11\\
\hspace*{4em}pour i de 0 a m-1 :\\
\hspace*{6em}z $\leftarrow$ 0\\
\hspace*{6em}pour j de 0 a 10 : z $\leftarrow$ z + th[j] * X[i][j]\\
\hspace*{6em}J $\leftarrow$ J + softplus(z) - y[i] * z\\
\hspace*{6em}e $\leftarrow$ sigmoide(z) - y[i]\\
\hspace*{6em}pour j de 0 a 10 : g[j] $\leftarrow$ g[j] + e * X[i][j]\\
\hspace*{4em}J $\leftarrow$ J / m\\
\hspace*{4em}pour j de 0 a 10 : g[j] $\leftarrow$ g[j] / m\\
\hspace*{4em}si J\_prec defini et abs(J\_prec - J) / max(1, abs(J)) < EPS :\\
\hspace*{6em}sortir de la boucle\\
\hspace*{4em}J\_prec $\leftarrow$ J\\
\hspace*{4em}pour j de 0 a 10 : th[j] $\leftarrow$ th[j] - ALPHA * g[j]\\
\hspace*{2em}theta[k] $\leftarrow$ th\\
\\
ecrire\_poids("weights.csv", VARIABLES, med, mu, sd, MAISONS, theta)
\end{pseudo}

La boucle sur $j$ dans le calcul de \texttt{z} et de \texttt{g} correspond au
produit $\X\thv_k$ et au produit $\X^\top(\mathbf p-\y_k)$. Une bibliothèque de
calcul matriciel remplace les deux boucles internes par
\texttt{z = X @ th} et \texttt{g = X.T @ (p - y) / m}, avec le même résultat.

L'ordre des opérations compte sur deux points. La division de \texttt{g} par $m$
arrive avant la mise à jour, et \texttt{th[j]} est modifié après le calcul
complet de \texttt{g} : les onze composantes descendent depuis le même vecteur.

Exécuté sur les $1\,600$ lignes, ce programme consomme $436$ itérations pour
Hufflepuff, $460$ pour Ravenclaw, $836$ pour Slytherin et $1\,069$ pour
Gryffindor, sous le plafond de $6\,000$. Les coûts finaux vont de $0{,}044$ à
$0{,}069$, et le modèle reclasse correctement $0{,}982$ de son propre ensemble
d'apprentissage.

\section{Le programme de prédiction}

La prédiction lit les paramètres, construit $\X$ avec les statistiques du fichier
de poids, calcule quatre scores par ligne et retient le plus grand, selon
l'équation~\eqref{eq:ovr-decision}. La sigmoïde est croissante et commune aux
quatre modèles : comparer les $z_k$ donne le même classement que comparer les
$q_k$.

La remarque s'étend à \rsym{softmax}{softmax}, dont le dénominateur est commun aux
quatre classes : la section~\ref{sec:lois} en donne l'argument. Les trois
quantités $z_k$, $\sigma(z_k)$ et $\mathrm{softmax}(z)_k$ désignent le même
argmax, pour tout jeu de scores. Sur les $1\,600$ lignes du fichier
d'apprentissage, les trois règles produisent la même maison, avec zéro ex æquo.

Deux mesures situent la difficulté du problème. L'écart entre le meilleur score
et le second vaut $2{,}05$ au minimum et $8{,}69$ en médiane. Le maximum des
quatre scores est positif sur chaque ligne, et une seule classe y dépasse $2$.
Les quatre maisons occupent donc des régions franchement séparées, ce qui éclaire
le résultat suivant : une régression softmax multinomiale, ajustée conjointement
sur les mêmes découpages, produit les $322$ mêmes prédictions que OvR à chaque
graine, pour une exactitude moyenne identique de $0{,}983$. Sur ce jeu de
données, le choix entre les deux modèles porte sur l'interprétation des sorties,
traitée au chapitre~\ref{ch:ovr}.

\begin{procedure}{logreg\_predict}
\textbf{Entrées :} \texttt{dataset\_test.csv} et le fichier de poids.
\textbf{Sortie :} \texttt{houses.csv}.
\begin{enumerate}
  \item lire l'ordre des variables, les statistiques de préparation, l'ordre des
        classes et les quatre $\thv_k$ ;
  \item construire $\X\in\R^{400\times11}$ avec ces statistiques ;
  \item calculer $\mathbf Z=\X\Theta^\top\in\R^{400\times4}$, où $\Theta$ empile
        les quatre vecteurs de paramètres ;
  \item pour chaque ligne, prendre l'indice du maximum et le convertir en
        libellé ;
  \item écrire l'en-tête puis $400$ lignes, dans l'ordre du fichier d'entrée.
\end{enumerate}
\end{procedure}

\begin{pseudo}[logreg\_predict dataset\_test.csv weights.csv]
VARIABLES, med, mu, sd, MAISONS, theta $\leftarrow$ lire\_poids(argv[2])\\
lignes $\leftarrow$ lire\_csv(argv[1])\\
\\
pour i dans lignes :\\
\hspace*{2em}si lignes[i][Astronomy] est vide et lignes[i][DADA] est rempli :\\
\hspace*{4em}lignes[i][Astronomy] $\leftarrow$ -100 * lignes[i][DADA]\\
\hspace*{2em}pour f dans VARIABLES :\\
\hspace*{4em}si lignes[i][f] est vide : lignes[i][f] $\leftarrow$ med[f]\\
\\
ouvrir sortie "houses.csv"\\
ecrire "Index,Hogwarts House"\\
\\
pour i dans lignes :\\
\hspace*{2em}x $\leftarrow$ [1.0]\\
\hspace*{2em}pour f dans VARIABLES :\\
\hspace*{4em}ajouter (lignes[i][f] - mu[f]) / sd[f] a x\\
\hspace*{2em}meilleur $\leftarrow$ 0 ; z\_max $\leftarrow$ -infini\\
\hspace*{2em}pour k de 0 a 3 :\\
\hspace*{4em}z $\leftarrow$ 0\\
\hspace*{4em}pour j de 0 a 10 : z $\leftarrow$ z + theta[k][j] * x[j]\\
\hspace*{4em}si z > z\_max : z\_max $\leftarrow$ z ; meilleur $\leftarrow$ k\\
\hspace*{2em}ecrire lignes[i][Index] + "," + MAISONS[meilleur]
\end{pseudo}

\begin{fichier}[houses.csv]
Index,Hogwarts House\\
0,Gryffindor\\
1,Hufflepuff\\
2,Ravenclaw\\
3,Hufflepuff\\
{[}...{]}\\
399,Gryffindor
\end{fichier}

La boucle sur \texttt{k} réalise l'argmax de l'équation~\eqref{eq:ovr-decision}
en une passe, sans stocker les quatre scores. L'initialisation de \texttt{z\_max}
à $-\infty$ garantit qu'une maison est retenue dès le premier tour, y compris
quand les quatre scores sont négatifs.

Les lignes de \texttt{dataset\_test.csv} dont plusieurs notes manquent reçoivent
une maison comme les autres : leurs cases vides ont été remplies par les médianes
d'apprentissage à l'étape~2. Le fichier produit compte $401$ lignes.

\section{Contrôles}

La colonne \texttt{Hogwarts House} de \texttt{dataset\_test.csv} est vide.
L'exactitude se mesure donc sur un ensemble de test découpé dans
\texttt{dataset\_train.csv}, selon le protocole du chapitre~\ref{ch:evaluation}.
Le sujet fixe le seuil d'acceptation à $0{,}98$, mesuré par l'évaluateur sur
\texttt{dataset\_test.csv}.

\begin{table}[htbp]
\centering
\begin{tabularx}{\textwidth}{X r}
\toprule
contrôle & valeur attendue \\
\midrule
gradient analytique contre différence finie centrée & écart relatif $<10^{-6}$ \\
sigmoïde et perte évaluées en $z=\pm1000$ & valeurs finies \\
exactitude d'une prédiction constante sur Hufflepuff & $0{,}331$ \\
exactitude après \rsym{permutation}{permutation} aléatoire des étiquettes & $\approx0{,}34$ \\
exactitude sur le test interne & $\geq0{,}98$ \\
lignes de \texttt{houses.csv} & $401$ \\
\bottomrule
\end{tabularx}
\caption{Valeurs de référence mesurées sur ce jeu de données.}
\label{tab:checks}
\end{table}

Le contrôle par permutation mélange les maisons entre les lignes avant
l'entraînement. Le chapitre~\ref{ch:evaluation} situe le résultat attendu à
$0{,}25$ pour quatre classes équilibrées. Les effectifs valent ici $529$, $443$,
$327$ et $301$ : privé de tout lien entre notes et maison, le modèle converge
vers la prédiction constante sur la classe majoritaire et atteint $0{,}34$. La
référence d'un contrôle par permutation est la performance de la meilleure règle
constante.

\begin{resultat}[Sortie de l'entraînement]
Quatre vecteurs de onze coefficients, dix médianes, dix moyennes, dix
écarts-types, la liste ordonnée des variables et celle des maisons. Ces $74$
nombres et ces deux listes forment l'intégralité de ce que lit la prédiction.
\end{resultat}
